\documentclass[11pt]{article}
\usepackage[a4paper,margin=28mm]{geometry}
\usepackage{amsmath,amssymb,amsthm}
\usepackage[colorlinks=true,linkcolor=blue,urlcolor=blue,citecolor=blue]{hyperref}
\newtheorem{theorem}{Theorem}
\newtheorem{lemma}{Lemma}
\newcommand{\coeff}[2]{\left[#1\right]#2}
\newcommand{\D}{\mathcal D}
\title{An integer proof of the Costa--Pellegrini coefficient identity}
\author{A Codex-assisted research project\\
  \small Project account: \texttt{xhao0307}}
\date{10 October 2026}
\begin{document}
\maketitle
\begin{abstract}
We prove, for every $k\ge2$ and every $1\le j\le k$, the equality between
two coefficient families discussed by Costa and Pellegrini in 2020.
The proof uses the signed expansion of a Vandermonde polynomial and an
involution of intervals. All steps are identities over the integers;
no interpolation denominators or coefficient nonvanishing assumptions
are required. The accompanying Lean project proves the full universal
statement. This identity does not settle Graham's rearrangement conjecture.
Publication priority and prize eligibility are not asserted.
\end{abstract}

\section{The published question}
Work in $\mathbb Z[x_1,\ldots,x_k]$, and define
\[
 T=\sum_i x_i,\qquad V=\prod_{i<r}(x_r-x_i),\qquad
 H=\prod_{1\le i<r\le k}(x_i+\cdots+x_r).
\]
Let $h$ be the product defining $H$ with its single full-interval factor
$T$ deleted, so $H=Th$. In particular $h$ is a polynomial before any
specialization. Put
\[
 c_{k,j}=\coeff{x_j^{k-2}\prod_{i\ne j}x_i^{k-1}}{Vh},\qquad
 e_{k,j}=\coeff{\prod_{i\ne j}x_i^k}{VH}.
\]
The distinguished variable has exponent zero in the second extraction.
These are the families in \cite[Section 4.1, p.~485]{CP}, where an
all-$k$ proof of their observed equality is requested.

\begin{theorem}\label{thm:main}
For every $k\ge2$ and $1\le j\le k$,
\[
 e_{k,j}=(-1)^{\binom{k-1}{2}}c_{k,j}
        =(-1)^{\lfloor(k-1)/2\rfloor}c_{k,j}.
\]
\end{theorem}

\section{Integer coefficient operations}
Set $M=\prod_i x_i^{k-1}$ and $\D(P)=\coeff{M}{VP}$.
Write $Z_j$ for setting $x_j$ to zero. A variable permutation $\sigma$
acts by substitution $P\circ\sigma=P(x_{\sigma(1)},\ldots,x_{\sigma(k)})$.

\begin{lemma}\label{lem:trace}
For every polynomial $Q$ and permutation $\sigma$,
\[
 \D(TQ)=0,\qquad
 \D(P\circ\sigma)=\operatorname{sgn}(\sigma)\D(P),\qquad
 c_{k,j}=\D(x_jh).
\]
\end{lemma}
\begin{proof}
The last equality is coefficient shifting. The middle equality follows
because $M$ is invariant under permutations and $V$ changes by their sign.
For the first equality it suffices, by linearity, to take
$Q=\prod_t x_t^{a_t}$. Expand
$V=\sum_r\operatorname{sgn}(r)\prod_t x_t^{r_t}$, where $r$ ranges over
permutations of the exponent labels $0,\ldots,k-1$.
A contribution to $\D(TQ)$ selects an index $i$ and satisfies
$r_t+a_t+\delta_{t,i}=k-1$ for every $t$.
Thus $r_i\le k-2$. Let $l$ be the unique index with $r_l=r_i+1$.
Swap the exponent labels at $i,l$ and select $l$ instead of $i$.
All resulting exponents remain unchanged, while the sign reverses.
This is a fixed-point-free involution of the contributing terms, so they
cancel over $\mathbb Z$.
\end{proof}

\begin{lemma}\label{lem:zero}
If $P$ is independent of $x_j$, then
\[
 \coeff{\prod_{i\ne j}x_i^k}{V(Z_jx)P}
     =(-1)^{k-1}\D(P).
\]
Consequently $e_{k,j}=(-1)^{k-1}\D(H(Z_jx))$.
\end{lemma}
\begin{proof}
In $\D(P)$, independence forces the Vandermonde exponent at $j$ to be
$k-1$. In $V(Z_jx)$, only terms with exponent zero there survive.
Cycle the exponent labels $0\mapsto1\mapsto\cdots\mapsto k-1\mapsto0$.
This bijects the two term families and multiplies the sign by $(-1)^{k-1}$.
At every other coordinate the Vandermonde exponent increases by one,
so the exponent required from $P$ is unchanged:
$k-(r_i+1)=k-1-r_i$. Apply this with $P=H(Z_jx)$; the exponent-zero
coefficient is unchanged by setting $x_j$ to zero.
\end{proof}

\section{A direct interval involution}
Fix $j$ and let
\[
 \rho_jx=(x_{j-1},\ldots,x_1,x_j,x_k,\ldots,x_{j+1}).
\]
Its sign is
$\alpha=(-1)^{\binom{j-1}{2}+\binom{k-j}{2}}$.
Let $\beta=(-1)^{k-1+(j-1)(k-j)}$.

\begin{lemma}\label{lem:boundary}
The integer polynomial $H(Z_j\rho_jx)-\beta x_jh$ is divisible by $T$.
\end{lemma}
\begin{proof}
Work in the quotient by $(T)$, a commutative ring; no domain cancellation
is needed. Use zero-based indices $0,\ldots,k-1$ in this paragraph,
with distinguished index $d=j-1$. Put $y=Z_d\rho_dx$.
For an inclusive interval $[a,b]$ its sum transforms as follows:
\[
 \sum_{i=a}^b y_i=
 \begin{cases}
  \sum_{i=d-1-b}^{d-1-a}x_i,& b<d,\\
  \sum_{i=k+d-b}^{k+d-a}x_i,& a>d,\\
  -\sum_{i=d-a}^{k+d-b-1}x_i,&a\le d\le b.
 \end{cases}
\]
In the crossing case the reflected left and right tails, together with
the displayed complementary interval, partition all original indices;
the total sum is zero in this quotient.
These formulas define an involution on all nonempty intervals. It swaps
the full interval $[0,k-1]$ and the singleton $[d,d]$; other singletons
remain singletons and all other strict intervals remain strict intervals.
Therefore the strict interval factors in $H(y)$ correspond bijectively
to the factors of $h(x)$ together with $x_d$.
The number of crossing strict intervals is
\[
 (d+1)(k-d)-1=k-1+d(k-1-d).
\]
Only these factors contribute a minus sign. Their product proves
$H(y)=\beta x_dh$ in the quotient, which is exactly the claimed
divisibility over $\mathbb Z$.
\end{proof}

\begin{proof}[Proof of Theorem~\ref{thm:main}]
Let $P=H(Z_jx)$. Since $\rho_j$ fixes $j$, it commutes with $Z_j$.
Lemmas~\ref{lem:trace} and \ref{lem:boundary} give
\[
 \alpha\D(P)=\D(H(Z_j\rho_jx))=\beta\D(x_jh)=\beta c_{k,j}.
\]
As $\alpha^2=1$, Lemma~\ref{lem:zero} yields
$e_{k,j}=(-1)^{k-1}\alpha\beta c_{k,j}$.
Put $u=j-1$, $v=k-j$. Then $u+v=k-1$ and
$\binom{u+v}{2}=\binom u2+\binom v2+uv$,
so the total sign is $(-1)^{\binom{k-1}{2}}$.
For $n=2r$ or $2r+1$, $\binom n2$ has the parity of $r$,
giving the other stated form.
\end{proof}

\section{Formal verification and scope}
The accompanying project uses Lean 4.34.1 and Mathlib revision
\texttt{d13f23b723b8a846827a245b89c10fc7d3f11612}.
Its theorem \texttt{CostaPellegrini.coefficient\_identity} quantifies over
all $k\ge2$ and all $j:\mathrm{Fin}\,k$, with precisely the above integer
coefficient definitions. The proof uses no numerical rank restriction,
coefficient nonvanishing premise or additional mathematical axiom.
The transitive axiom audit reports only
\texttt{propext}, \texttt{Classical.choice}, and \texttt{Quot.sound}.

This theorem equates two coefficient vectors. It does not prove that
either vector is nonzero modulo every relevant prime and therefore does
not settle Graham's original conjecture. The question's published source
is identified here, but worldwide priority, independent human peer review,
catalogue inclusion and reward entitlement have not been established.
The project was initiated by its human requester; Codex developed the
mathematical arguments and Lean scripts with coordinated AI reviews.

\begin{thebibliography}{9}
\bibitem{CP} S. Costa and M. A. Pellegrini,
\emph{Some new results about a conjecture by Brian Alspach},
Archiv der Mathematik \textbf{115} (2020), 479--488.
\href{https://doi.org/10.1007/s00013-020-01507-7}{doi:10.1007/s00013-020-01507-7}.
\end{thebibliography}
\end{document}
